\chapter[Non-stationary Electrochemical Impedance spectroscopy]{Non-stationary Electrochemical Impedance spectroscopy: state of the art}
% \minitoc

\section{Historical development}

\subsection{Single sine impedance}
We need to start from the basics of the practical measurement of Electrochemical Impedance Spectroscopy to understand the use of multi-frequency waveform and measurements in non-stationary conditions. During the 1970s impedance or admittance was already measured for electrochemical systems over a wide range of frequencies for the study of the kinetics of charge transfer and transport after the pioneering works of Randles (1947) based on Warburg formalization of transport impedance in 1889. The practical methods to measure impedance was though alternate current bridges methods (where an unknown impedance elemtn is put in a network of known resistors with a certain design) and using a lock-in amplifier with a sinus function generator. The lock-in amplifier is an analog electrical component capable of extracting the amplitude and phases of a signal frequency component from a noisy signals. The frequency is dictated by the choice of the sinus waves. The frequency was manually changes and a full spectra could take a time of a few minutes per decade. At that time the frequency band was limited to 100kHz -10mHz for the hardware. The lock-in amplifiers technologies was quite mature at that time starting from the commercialization from Princeton Applied  Research (model 124A) by Robert H. Dicke. 
[Troppi dettagli e informazioni discordandi da LLM. Meglio dire solo del modello 1205. Il 1174 usava dell'elettronics per automatizzare lo sweep e data-handling] The breakthrough became with the production of the Solartron 1174 in 1978 from Princeton Applied  Research. It integrates for the first time a frequency response analyzer (FRA) to compute the impedance using digital cross-correlation. The frequency band was also extended to the range 10 $\mu$ Hz and 1 MHz. With the following model 1250 in 1983 potential controll, automatic frequency sweep and digital signal processing brings the de facto standard of electrochemical impedance measurement. The measurement time is reduces to 10-30 minutes. 
Modern workstations still use the same concept. Circuits are now more complex, though, allowing to apply a bias voltage or current. They also use avaraging and automaticocurrante range selection to increase signal to noise ratio over all decades.
Nowadays FPGA offer fast, inexpensive and reliable digital signal processing. Their are the core of modern electrochemical workstation as well as audio processing and any other signal application.

\subsection{Multi-frequency excitation}
During the 1970 the revolution of the microcomputer was also taking place allowing the first digital data processing. Smith and Creason in 1972 dintegrated an algorithm into a micrcoputer for the online computation of admittance based on the Fast Fourier Transform 

In the years 1971, 1972 and 1973 a seires of pubblications showed the possiblity of obtaing the impedance using algorithms based on the Fast Fourier Trasnform performed by microcomputers. Kojima and Fujiwara \cite{kojima1971fourier} and later Glover and Smith \cite{gloverAlternatingCurrentPolarography1973} showed the possibility of perfoming "alternated polarography" with digital signal processing on microcomputers. Creason and Smith applied the method for the estimation of faradaic admittance in fast reactions \cite{creasonFourierTransformFaradaic1972} and performed a screening of possible multif-requency exitations \cite{creasonFourierTransformFaradaic1973a}.

The method is much faster than using the single sine approach and resulted in a better signal to noise ratio thanks to the averaging across frequencies. 
Despite being promising, the microcomputers were still expensive and the generation of the multi-frequency perturbation require linear amplifiers with large dynamic range and precise and phase-coherent waveform generators. The use of multi-frequency waveforms is prone to rise non-linearities that generates intermodulations. The single sine method proves to be much simpler and good enough in terms of signal to noise ratio to became the de facto standard in electrochemistry teaching and research slowing interrupting the development of the multi-sine method even with the increase of hardware capabilities at low price. The method was fast as 1s but the frequency range was between around 100Hz and 40 kHZ.

\subsection{Impedance of evolving systems}
Bond, Schwall and Smith in 1977 \cite{schwallOnlineFastFourier1977} combined the method of the online FFT using cyclic voltammetry and multi-frequency waveform testing and applied it on mercury and platinum electrodes \cite{bondOnlineFFTFaradaic1977} \cite{schwallOnlineFastFourier1977}. The first dynamic impednace measuremnt is born. They were also working on the optimization of the implementation on the microcomputer arhitecutre \cite{schwallHighSpeedSynchronous1977}. 
In the following years, Smith collaborated with many reseracher to apply the method, noe baptized "fourier transform faradaic admittance measurements (FT-FAM)" for the investigation of electron transfer \cite{hayesACPolarographyUsing1980} \cite{grzeszczukMeasurementsMarcusTheory1983} \cite{grzeszczukStructuralConsequencesElectron1983} \cite{grzeszczukElectrochemicalInvestigationsTetraphenylethylene1984} \cite{ivaskaApplicationOnlineFast1985} \cite{grzeszczukElectronHydrogenAtom1986}. In 1990, Ivaska and collegues ported the technique into a smaller form factor with details on the hardware \cite{engblomExperiencesOnLineFourier1990}.

In 1996 Hazi, Bond and collegues, extended the system to modern hardare for the measurement of admittance or impedance in the frequency range 50Hz-50kHz during staircase voltammetry \cite{haziMicrocomputerbasedInstrumentationMultifrequency1997} \cite{schieweUnifiedApproachTrace1998}.

In 1982 T. Osaka and L. Naoi \cite{osakaApplicationOnlineImpedance1982} also developed their set-up using commercial interfaces to compute the FFT online for the estimation of the impedance fast and accuratly (Since I cannot get acces to the original publication I assume the frequency range is similar to the one proposed by Smith, onlyl they focus on the impedance rather then the admittance), that they used to study oxigen evolution \cite{osakaStudyTimedependenceOxygen1984} and processes in doped polymers \cite{osakaElectrochemicalStudyAnion1984} \cite{osakaStudyChargeDischargeCharacteristics1984}\cite{osakaApplicationConductive985}.

In 1986 T. Tsuru and S. Haruyama \cite{1986296} proposed a review paper on the differences between the 
highlighting the pros/cons and giving even a forcast of the future development and that despite the .... 
[FRA systems excel through their methodical frequency-by-frequency measurement approach, delivering superior accuracy and exceptional noise immunity that allows meaningful results even when signals are completely buried in noise. This systematic approach makes them ideal for high-precision measurements of stable systems, though at the cost of longer measurement times since each frequency requires at least one complete cycle. Conversely, FFT systems offer dramatic speed advantages by simultaneously analyzing multiple frequencies, theoretically measuring wide ranges like 500Hz to 1Hz in approximately one second, while providing versatility in signal processing with various waveforms including pulses, steps, and white noise. However, FFT systems traditionally suffered from lower measurement accuracy, greater susceptibility to external noise, and frequency resolution bias toward high-frequency regions, requiring more sophisticated signal conditioning and averaging operations. They envisioned FRA systems focusing on impedance and harmonic measurements with computer-controlled frequency selection and display functions, while FFT would develop into versatile tools for analyzing electrochemical transient characteristics from multiple perspectives. ]
[While integrated circuit advances have significantly modified these fundamental trade-offs, core differences persist in modern implementations. Contemporary digital signal processors and FPGAs have dramatically improved FFT accuracy and noise handling, narrowing the performance gap that once strongly favored FRA systems, while reducing cost differentials between approaches. Modern instruments increasingly employ hybrid strategies, automatically selecting optimal measurement methods based on specific requirements - leveraging FRA-type approaches for high-precision single-frequency measurements and FFT-type methods for broadband analysis and rapid transient monitoring. However, the basic relationship between measurement time and accuracy per frequency point continues to favor methodical approaches for precision work, while simultaneous multi-frequency analysis remains advantageous for speed-critical applications. Today's impedance analyzers often embody the authors' vision of intelligent, software-defined instruments that synthesize both approaches rather than forcing users to choose between competing methodologies, though the fundamental trade-offs between systematic precision and simultaneous analysis speed remain relevant considerations in modern electrochemical research.]

\subsection{The first time-resolved impedance spectrometer}
\label{sec:fft-eis}
The modern set-up for FFT based impedance estimation was presented by G.S. Popkirov and R.N Schindler in 1992 made of the combination of home made potentiostat, a dedicated waveform generator and an acquisition device (called in the text "transient recorder").The potentiostat has two types of amplifier to work in the frequency range 100kHz-100Hz or 0.76Hz-0.76mHz. The autor tested firrent types of perturbations: rectangular pulse, randomly spread rectangular pulse, pseudorandom noise, serial sequence of sinusoids, consequent sinusoidal with decreasing frequency, and a waveform of summed 42 sinusoidal of spread frequencies across the range. The latter gave the best signal to noise ratio on the desired frequency and large band. The authors also discuss the optimization of phases and amplitude of the sinusodial components of a multi-frequency waveform \cite{popkirovOptimizationPerturbationSignal1993a} and the validation of the measurement for the interference of non-linearities \cite{popkirovValidationExperimentalData1993a} \cite{popkirovEffectSampleNonlinearity1995a}.
Throught the 90s Popkirov and Schindler applied the method for electropolimerization \cite{popkirovInSituTimeResolved1993}, the characterization of redox process in conductive polymers \cite{popkirovInsituImpedanceMeasurements1995} \cite{popkirovElectrochemicalImpedanceSpectroscopy1995} \cite{popkirovInvestigationConductingPolymer1997} \cite{garcia-belmonteDeterminationElectronicConductivity2003a}, silver dissolution \cite{popkirovElectrodePotentialRedistribution1995} \cite{baltrunasTimeEvolutionActive1997} and Benzophenone/Benzophenone anion radical couple in different solutions \cite{walczykInvestigationRedoxCouple1995}.
They used both sweeping voltage and galavnostatic techniques with a multi-sine that covers four and half decades from 1 Hz to 31 kHz.
I calcoli della FFT sono fatti da un microcomputer o normale PC che riceve i dati?

\subsection{Short-time Fourier transform (STFT) viewpoint}

Identification of slowly-varying dynamics is analogous to estimating the spectrogram of a signal:
\[
Y(t,\omega) = \int_{-\infty}^{\infty} y(\tau)\, w(\tau - t)\, e^{-j\omega \tau} d\tau,
\]
where $w$ is a window of width $\Delta t$. The corresponding transfer function estimate is:
\[
G(t,\omega) \approx \frac{Y(t,\omega)}{U(t,\omega)}.
\]
This representation makes explicit the need to choose a compromise between:
\begin{itemize}
\item \textbf{time resolution} (small windows, good tracking of fast parameter variations), and
\item \textbf{frequency resolution} (large windows, accurate FRF estimation).
\end{itemize}

This trade-off is fundamental and cannot be eliminated; it mirrors the time-frequency uncertainty principle.

\section{Limitations of LTV identification and fundamental trade-offs}

\subsection{Resolution trade-off: the “Heisenberg inequality” of system identification}

For a window of effective width $\Delta t$, the variance of a frequency estimate decreases with $\Delta t$, while the bias due to time-variation increases with $\Delta t$. A time-varying FRF can only be estimated accurately if:
\[
T_{\mathrm{var}} \gg \Delta t \gg T_{\mathrm{dyn}},
\]
where $T_{\mathrm{var}}$ is the characteristic parameter drift time. This inequality provides a design rule for multisine and broadband excitations.


\subsection{Modern set-ups}

From 2000 to 2005 D. Roy, J. E. Garland, C.M. Pettit and collegues from the department of physics at Clarkson University in New York implemented a set up  similar to Popkirov's design but using commercial device controlled via Lab View in a personal computer \cite{waltersWeakAdsorptionAnions2001a}. The multi-sine perturbation has 133 components with odd frequencies, random phase and equal amplitude. The estimation of the impedance is via Short time fouerire transform on a segment of unwindowed signals. The discussess also in details how to choose parameters and the experimental limitation of such tehcnique \cite{garlandAnalysisExperimentalConstraints2004}. They appliead the technique to study difference processes on gold such as adsorbtion of ions, deposition of bismuth \cite{garlandKineticAnalysisElectrosorption2002a} and indium with focus on coadsorbtion, electrooxidation of methanol \cite{assiongbonElectrooxidationMethanolGold2005}, redox couple kinetics with focus on coadsorbtion \cite{emeryTimeResolvedImpedance2005} and chemical reduction of Nickel \cite{shiElectrocatalysisHypophosphiteNickel2014}. 
They also coupled time-resolved impedance with optical measurements during indium deposition on molybdenum \cite{pettitElectrodepositionIndiumMolybdenum2002}. Furthermore they studied surface reaction and effect planarization for different methals, Tallium \cite{assiongbonChemicalRolesPeroxidebased2004} \cite{gorantlaChemicalEffectsChemical2005} \cite{pettitRoleIodateIons2005}, copper \cite{luRelativeRoles22004} \cite{pettitTimeResolvedStudy2004} \cite{gorantlaCitricAcidComplexing2005} and silver \cite{assiongbonElectrooxidationMethanolGold2005} .

From 2003 on, K. Darowicki, P. S'lepski and their collegues, they put together different instrumental set-ups that follow the schematics from Popkirov of waveform generator, potentiostat and data sampler, but using always commercially available device controlled by computer \cite{darowickiInvestigationsPassiveLayer2003} \cite{darowickiInstantaneousElectrochemicalImpedance2004a}. The impedance is still computed using the short-time Fourier Trasnform, at first imployin exponentail chirp perturbations (in the range 1320–144 Hz) \cite{darowickiDynamicElectrochemicalImpedance2003a} \cite{darowickiContinuousFrequencyMethodMeasurement2003} \cite{darowickiDeterminationElectrodeImpedance2004a} and later with multisine exitation \cite{darowickiInvestigationsPassiveLayer2003}. 
The prolific work conducted by the group touches multiple applications: corrosion evaluation on stainless steel \cite{darowickiEvaluationPittingCorrosion2004a} \cite{darowickiTimeDependencePit2004} \cite{krakowiakImpedanceInvestigationPassive2005} \cite{krakowiakImpedanceMetastablePitting2005} \cite{orlikowskiInstantaneousImpedanceMonitoring2015} \cite{rylInstantaneousImpedanceMonitoring2016} \cite{orlikowskiDeterminationPittingCorrosion2017} \cite{naroznyApplicationElectrochemicalImpedance2018} \cite{darowickiNovelApplicationDynamic2020a} \cite{rylInstantaneousImpedanceAnalysis2017} and metal alloys \cite{arutunowElectricalMechanicalCharacterization2015} \cite{zielinskiApplicationMultisineNanoscale2015} \cite{wysockaInvestigationElectrochemicalBehaviour2016} \cite{gawelImpedanceMonitoringCorrosion2017}, monitoring of the crack of corrosion layers on stainless steel \cite{darowickiDetectionStressCorrosion2004}, characterization of conductive polymers \cite{darowickiImpedanceCharacterizationProcess2004} \cite{darowickiDynamicElectrochemicalImpedance2004}, investigating the hydrogen storage alloy[2010], monitoring the discharge of nickel/nickel hydride batteries[2013],  monitoring the adhesion of organic coating to stainless steel and galvanic coupling \cite{burczykLocalElectrochemicalImpedance2018}.
The frequency range of the multi-sine of these works is around 3 Hz to 20 kHz. The method was also coupled with atomic force microscopy for the study of surface terminations of boron-doped diamond \cite{zielinskiMultifrequencyNanoscaleImpedance2019} \cite{rylStudySurfaceTermination2016} and the quartz crystal microbalance \cite{parasinskaApplicationQuartzCrystal2024}. 
The researchers also applied the techniques to Fuel cell under different operation \cite{slepskiImpedanceMonitoringFuel2015} \cite{darowickiInfluenceDynamicLoad2019} \cite{darowickiImpedanceHydrogenOxidation2020a} and Redox flow batteries \cite{krakowiakThreeModesElectrochemical2024}. Here the frequency range was extended to 4.5 kHz–0.3 Hz (or 30 mHz). 

In 2004 J. Schouckens and R. Pintelon started to apply their theoretical framework [???] of system identification in the frequency domain characterization of electrochemical systems with A. Hubin to obtain the electrochemical impedance of corroding metals using a set-up has the same strucutre of the one from Popkirov but usign multi-sine perurbation with low frequency components down to 0.01Hz \cite{vangheemElectrochemicalImpedanceSpectroscopy2004a}\cite{vangheemElectrochemicalImpedanceSpectroscopy2006a} \cite{blajievImprovementImpedanceMeasurement2006}. It is interesting how they compare multi-sine and stepped sine consecutive measurement to validate the experiment in \cite{vangheemInstantaneousImpedanceMeasurements2004}. In 2009 they discuss the advantage of using only the odd-number harmonics in the sequence of summed sinusoidal with 220 components between around 10mHz and 2kHz \cite{vaningelgemAdvantagesOddRandom2009}. 
In particular, they put attention on the estimation procedure and the calculation of the impedance variance and the error due to the non-linearities and non-stationarity of the system (quantify them). With their method, the istantanous impedance is estimated from consecutive periods of the mutlisine or continously in the same fashioin of the short time fourier transform
Prior \cite{schoukensBroadbandSteppedSine1999} they also discussed the difference in estimating the frequency response function with multi-frequency waveform and stepped sine trhough simulation adresing signal to noise ratio and measurement time. Previously \cite{schoukensSurveyExcitationSignals1988a} they study quality and length of different excitation signals  including periodic signals: stepped sine, swept singe (chirp), multi-sine, pseudorandom noise, periodic noise, maximum length binary sequence and discrete interval binary sequence ; transient signals : impulse, random burst; and aperiodic signals : random noise.
Their method was applied in multiple studies up to this day: corrosion of copper\cite{vaningelgemApplicationMultisineImpedance2008}, alluminum and stainles steel \cite{breugelmansOddRandomPhase2010a} \cite{alvarez-pampliegaCorrosionStudyAlrich2014} \cite{jiOddRandomPhase2018} \cite{jiElectrochemicalRamanAnalysis2020}, the adsorbtion of molecules on metals \cite{blajievPotentiodynamicEISInvestigation2008} \cite{geboesInvestigationAdsorptionMechanism2015} and oxides \cite{hauffmanMeasuringAdsorptionEthanol2012}, self-assembling monolayers \cite{hauffmanDynamicSituStudy2013}\cite{fernandezmaciaStudyElectronTransfer2014}, polymeric coatings \cite{jorcinInvestigationSelfhealingProperties2010}\cite{breugelmansOddRandomPhase2010}\cite{woutersMonitoringInitialContact2021} \cite{meeusenFEMModellingPredict2023}\cite{madelatSituMonitoringLateral2023}, ferri/ferrocyanide reaction on gold \cite{fernandezmaciaORPEISStudyTime2015}, aptomer as biosensors \cite{pauwelsNewMultisinebasedImpedimetric2016} and alloy protection \cite{meeusenUseOddRandom2018}. 

In 2009 R.L Sacci and D. A. Harrington from the Department of Chemistry of the University of Victoria in British Columbia (Canada) proposed their set-up following the schematics of Popkirov using commercial devices and multi-frequency excitation coupled with cyclic voltammetry. The frequency band of the multi-sine is 1Hz - 1.3 kHz with removed harmonics to avoid intermodulation following Popkirov 1993 and amplitude decreasing by a factor of two every two decades.
They used the method to study electroxidation of metanol  \cite{sacciDynamicImpedanceFormic2009}
\cite{holmUnderstandingReactionMechanisms2018} \cite{holmDynamicElectrochemicalImpedance2019} \cite{holmUnderstandingReactionMechanisms2019} ,  electrocatalytic reactions \cite{sacciDynamicElectrochemicalImpedance2014} and droplet interface bilayers \cite{sacciDisentanglingMemristiveMemcapacitive2022}.

in 2018 H. Zappen D. U. Sauer and collegues from RWTH Aachen Univeristy in Germany, used the usual schematic of Popkirov with only one FPGA based potentiostat board of which they could modify the software to generate a multi-frequency waveform with 20 sinusoids in the frequency band 1 Hz - 1kHz \cite{zappenApplicationTimeResolvedMultiSine2018}. In their work, they put enphasis on the design of the multisine and the screening of the main phase optimization strategies in both efficiency and effectivnes (resulting crest factor). The impedance is measured in equilibrium state and during externally induced thermal runaway (no drifts) coupling ultrasonic measurements \cite{zappenInOperandoImpedanceSpectroscopy2020a}. Zappen later openeed a startup, Safion GmbH, for the use of the method for fast quality assessment of battery after production with a dedicated device.

In 2020 T. Pajkossy and G. Mészáros proposed a set-up in which the potentiostat is custom build in order to analogically seperate the drift and oscillations. The input controll is still made by the superposition of a main drift, sweeping voltage, and small multi-frequency excitation with odd prime frequencies in the band 18.3 Hz - 17.8 kHz \cite{pajkossyConnectionCVsImpedance2020}. The four outputs are recorded with analog to digital converter. The researchers did also an extendive theoretical evaluation of the coupling of cyclic voltammetry and impedance for the estimation of charge transfer resistance for quasi-reversible redox systems \cite{pajkossyDynamicElectrochemicalImpedance2019}  \cite{pajkossyVoltammetryCoupledImpedance2020} \cite{pajkossyDynamicElectrochemicalImpedance2021}. 

\subsection{Non-stationary impedance}
The methods discussed so far in the section works for systems that can be approximated stationary for the period of the excitation. It is limiting the acquisition of impedance for fast systems or using high frequencies (usually above 1Hz). 
\label{sec:BLTVA}
From 2012 T. Bruegelmans, J. Latier and their colleagues in Brussels extended the theory of system identification for non-stationary systems to the estimation of electrochemical impedance \cite{breugelmansOddRandomPhase2012}  for system that are effected by a strong drift and being still able to estimate the impedance to frequency as low as 0.3 Hz. 
Later, the method has been further extended with the identifiation of non-linearities \cite{hallemansNonparametricIdentificationLinear2020} \cite{hallemansDetectionClassificationQuantification2021} and generalized for any non-statoinary systems \cite{hallemansBestLinearTimeVarying2021}. In Halleman 2022 they also introduce a new method for detrending the signal before calculate the best linear time-varying approximation.
From 2018 on, the group applied the method for the identification of the transfer function of many highly non-stationary systems performing galvanostatic techniques: electrograining \cite{woutersUseInstantaneousImpedance2018}, human respiration \cite{battistelHarmonicAnalysisSeparation2021},  metal electrodeposition \cite{colletOperandoORPEISStudy2021} \cite{dabirihavighOperandoOddRandom2022}, the conversion of zinc to zirconia \cite{dabirihavighOperandoOddRandom2023} and to study fast charging \cite{hallemansOperandoElectrochemicalImpedance2022}and SEI formation in Li-ion batteries \cite{dabirihavighOperandoORPEISMonitoring2025}.

\label{sec:dmfa}
In 2017 F. La Mantia, A. Battistel, D. Koster and collegues at the University of Bremen in Germany using the same set-up of Popkirov with a combination of commercial devices using a multi-sine superimposed with a cyclic voltammetry \cite{kosterDynamicImpedanceSpectroscopy2017} for fast measures of kinetic processes. The group doesn't use the short time fouerier trasnform but developed method for the time-frequency analyis of the highly non-linear and non-stationary multi-frequency signals filtering the data in the fourier space \cite{battistelPhysicalDefinitionDynamic2019}. This filtering behaves like a convolution of a sinc function of the original signal, a concept similar to the wavelet transform in signal processing but deduduced from intuition.
The group hilighted how the calculation of the time-resolved impedance with their method analytically converges to the short time fourier transform when the filter is as wide as the signal lenth. The researchers also discussed how the choice of the frequencies poresent in the multisine impact the quality of the impedance estimation with thier method and how to reduce the drift impact on the frequency sepctrogram a quasi-triangular wave was designed to replace the classic triangular voltage waveform used in the cyclic voltammetry. 
The method was used to study the redox active species \cite{kosterDynamicImpedanceSpectroscopy2017} [Mugisa 2024], disolution of p-type silicon \cite{kosterMeasurementAnalysisDynamic2018}, the intercalation of cations in Nickel Hexacyanoferrate \cite{erinmwingbovoDynamicImpedanceSpectroscopy2019}, lithium manganese oxide \cite{erinmwingbovoDynamicImpedanceSpectroscopy2020} [Erinwingbovo 2023] and lithium titanium phosphate \cite{scarpioniElectrochemicalImpedanceSpectroscopy2023a} and for the adsorbtion of ion on porous carbon electrodes. (The latter two are part of the discussion of this thesis).
The group also fronted the problem of fitting the great ammount of data from time-resolved impedance measurements \cite{battistelAnalysisNonstationaryImpedance2016} \cite{chukwuDynamicImpedanceSpectroscopy2023}. In 2025 the group published a work using piece-wise linearization approach from dynamic system theory to model the non-stationary impedance of an electrochemical system and validate the results of the Dynamic Multi-Frequency Analysis with bassband frequency filtering. Fruthermore the difference between stationary and non-stationary impedance is addressed. Is it correct ???

\subsection{Honorable mention}

Another important reserach to be mentioned is the work of Stoynov on the mathematical determination of the instantanous impedance from data sets of impedance subjected to non-stationariteis. He first tintroduced in 1985 the four-dimentional analysis for non-stationary systems [Stoynov 1985-1992]] and then in 1992 he introduce the rotating fourir transform [Stoynov 1992].
In this work, I address only the methods based on dynamic system identification with multi-frequency escitations so these two approaches, while theroetically interestng, will be not further discuss. 


\section{Nomenclature}

The method for identify the impedance of systems outside equilibrium, for being still in development, doesn't have a named agreed upon in the community. Let start from the single stepped sine impedance method based referred as classic, traditional, static or stationary. The immediate interposition to static, used in the sense of a system that is not changing (i.e stationary according to system theory jargon), is to call it dynamic. This word give an immediate picture in the mind of a system that changes with time. On the other hand, even the classic experiment is dynamic in nature in the sense that the system is drifted minimally and reversibly from its thermodynamic equilibrium state. The beginning and final state remain then unperturbed and we can use the definition of stationary. The type of measurement described in this work are instead applyed on systems that change with time, i.e non-stationary systems. Therefore the inequivocal name for the method should be non-stationary EIS to highlight the superposition of the characteristic time constant of the drift and excitation frequencies. The result of the experiment is the estimation of a time-varying impedance. I also use here the word estimation instead of calculation to highlight the presence of many assumption. The obtained frequency response function from the experiment might not be the real linear impedance response. That happen when the assumption are not verified in the practical system. For example the system varies to fast.

Time-resolved multi-sine electrochemical impedance spectroscopy is the best name for the technique. Impedance is the name of the measured quantity in any case. Time-resolved adjective can add information on the change to the impedance. Time-varying is a better adjectvie for the system.

Popkirov/schindler calls it also time domain eis.

Galvanodynamic? How can one also specificy the type of drift in the name of the technique?

NSEIS nonstationary electrochemical impedance spectroscopy


\section{Embedded specttrometers}
